% Original applied explanation; see Chapter04/SOURCES.md.
\subsubsection{গ্রন্থাগারের সূচি দিয়ে সার্চ ইঞ্জিন বোঝা}
গ্রন্থাগারের সূচিতে বইয়ের নাম ও অবস্থান লেখা থাকে; সূচি দেখলেই বইয়ের সব পৃষ্ঠা পড়া হয় না। সার্চ ইঞ্জিনও আগে খুঁজে পাওয়া পেজের তথ্য দিয়ে সূচি তৈরি করে। \textbf{কলেজের ভর্তি বিজ্ঞপ্তি} লিখলে সেই সূচি থেকে প্রাসঙ্গিক ফল দেখায়; ফলের লিংকে ক্লিক করলে ব্রাউজার আসল সাইটের পেজ খোলে।

কোনো পেজ নতুন হলে বা খোঁজার জন্য উন্মুক্ত না থাকলে তা সার্চ ফলের সূচিতে নাও থাকতে পারে। আবার সূচির তথ্য পুরোনো হতে পারে। তাই ফলের সারাংশের পাশাপাশি আসল পেজের প্রতিষ্ঠান, প্রকাশের তারিখ ও বিষয় মিলিয়ে দেখবে।
